% 开发者用大模型进行代码生成，依赖包选择
In recent years, large language models (LLMs) have developed rapidly and demonstrated strong capabilities across various software engineering tasks~\cite{wang2023review,jimenez2024swe,roziere2023code,sun2025source}. Developers increasingly rely on LLMs not only to generate code, but also to recommend libraries, configure dependencies, and prepare development environments~\cite{alhanahnah2024depsrag,latendresse2025robust}. In these workflows, package names generated by LLMs are often treated as actionable recommendations: they may be copied into import statements, installed through package managers, or used by automated coding agents. Consequently, the reliability of package-related outputs has become an important requirement for the practical adoption of LLMs in software development~\cite{pearce2025asleep,perry2023users}.



However, recent studies have shown that LLMs may suffer from hallucinations, generating outputs that appear plausible but are factually incorrect~\cite{zhang2025siren,gao2025current,zhang2025llm}.
One particularly concerning type is \emph{package hallucination}, where models generate non-existent or invalid software package names~\cite{spracklen2025we,krishna2025importing,zhao2025hfuzzer}. Unlike many conventional code generation errors, package hallucinations can directly affect dependency management and software supply chain security~\cite{ohm2020backstabber,ladisa2023sok,neupane2023beyond}. Developers and automated agents may inadvertently install hallucinated packages generated by LLMs, creating opportunities for attackers to register malicious packages under the same names~\cite{spracklen2025we}. Such risks are closely related to package confusion attacks, a well-known threat to software supply chains~\cite{lanyado2024diving,kaplan2021survey}. As LLM-based coding assistants and autonomous agents are increasingly integrated into development workflows, mitigating package hallucination is essential for building trustworthy software engineering environments.

Existing studies have proposed various strategies to mitigate package hallucinations, ranging from external mitigation techniques to internal model adaptation approaches~\cite{spracklen2025we}. External mitigation methods primarily operate outside the model. One line of work augments generation with external knowledge sources (e.g., PyPI~\cite{pypi}) to ground model outputs in reliable package information~\cite{lewis2020retrieval,jain2025mitigating}. Another line of work encourages the model to refine its own responses through self-correction~\cite{ji2023towards,madaan2023self}. Although these methods can reduce hallucinations, they do not directly modify the model’s internal behavior and introduce additional inference overhead or pipeline complexity~\cite{huang2025survey}.
Internal adaptation methods (e.g., supervised fine-tuning) provide a more direct alternative by modifying the model itself~\cite{spracklen2025we,tonmoy2024comprehensive}. However, fine-tuning updates a large number of parameters and may overfit to specific output formats, making it computationally expensive and potentially degrading the model’s general capabilities. These limitations motivate the need for a lightweight approach that can directly mitigate package hallucinations in LLMs without extensive retraining or architectural modifications.


Recently, model editing, also known as knowledge editing, has been proposed to update specific knowledge or behaviors in LLMs without full retraining~\cite{li2025model,wang2024detoxifying,zhang2024comprehensive,liu2025creme}. 
This provides a promising direction for mitigating package hallucination.
% , since the goal is to correct harmful package-related behavior while preserving the model’s general capabilities.
However, recent studies have shown that existing editing methods may still struggle to correct hallucinated outputs~\cite{huang2025can}.
Existing model editing methods primarily focus on factual knowledge updates, such as modifying subject-object associations or correcting individual facts~\cite{dai2022knowledge,li2024pmet,meng2022rome,wu2023depn}. 
These methods often rely on clear subject tokens or target phrases to locate and edit the corresponding knowledge. Package hallucination presents a different challenge. Package-related prompts usually contain complex task descriptions rather than explicit factual subjects, and the desired behavior is not to generate a fixed target output. Instead, the model should suppress non-existent package names while preserving valid packages. Furthermore, valid and hallucinated packages may appear together in the same response. This suggests that package hallucination is better viewed as a package-validity boundary problem than a factual rewriting problem. We define the package-validity boundary as the model's ability to distinguish valid packages (i.e., packages that exist in the package ecosystem) from hallucinated packages under a given task context, which is represented by the dashed separation line in Figure~\ref{fig:motivation}.

In this paper, we introduce \textbf{\appname}, a lightweight localized model editing framework for mitigating package hallucination in LLMs. Unlike existing mitigation methods that rely on external package retrieval or full-model fine-tuning, \appname aims to refine the package-validity boundary within the model itself. 
Specifically, \appname first employs a risk-aware localization strategy to identify modules whose parameters are most sensitive to the model’s preference for hallucinated over valid packages, which serves as a signal of the package-validity boundary. \appname then injects lightweight LoRA adapters into the localized modules and optimizes them with a boundary-aware objective that reinforces valid packages, suppresses hallucinated packages, and preserves the original model behavior on locality prompts.


To evaluate the effectiveness of \appname, we conduct experiments on a package hallucination dataset derived from package-related prompts~\cite{spracklen2025we}. We evaluate three representative open-source LLMs: \texttt{DeepSeekCoder}, \texttt{Qwen3}, and \texttt{Llama-3.1}. To provide a comprehensive evaluation, we compare \appname with five baselines, including two package hallucination mitigation methods (i.e., Full-FT~\cite{spracklen2025we} and Self-Refinement~\cite{spracklen2025we}) and three model editing methods (i.e., ROME~\cite{meng2022rome}, MEMIT~\cite{meng2023memit}, and DINM~\cite{wang2024detoxifying}).
Experimental results show that \ding{182} \appname effectively mitigates package hallucinations in the package recommendation task, reducing package-level hallucination rate (Package-HR) by 79.9\% on the edit set, and by 65.4\% on unseen prompts. \ding{183} The learned package-validity boundary generalizes to other package-related tasks, reducing Package-HR by 12.8\% in code generation and by 34.0\% in \texttt{pip install} recommendation. \ding{184} Compared with existing baselines, \appname achieves a better balance between hallucination reduction and valid package preservation, while several baselines reduce hallucinations by suppressing package generation. 
\ding{185} Ablation results show that both module localization and the three editing objectives are essential to effective editing.  
\ding{186} \appname is lightweight in practice, requiring only 731 KB of LoRA parameters on average, approximately 20,000$\times$ smaller than full fine-tuning checkpoints.

\noindent\textbf{Contributions:} The main contributions of this paper can be summarized as follows:
\begin{itemize}[leftmargin=*]
    \item We formulate package hallucination as a \emph{package-validity boundary} editing problem, aiming to refine the model's boundary between valid and hallucinated packages.
    \item We introduce \textbf{\appname}, a lightweight localized model editing framework that identifies modules related to package hallucination and applies boundary-aware LoRA editing to reduce hallucinated packages while preserving valid package recommendations.
    \item We conduct experiments on three open-source LLMs and show that \appname effectively reduces package hallucination, generalizes to other package-related tasks, and requires only lightweight editing artifacts.
\end{itemize}

